\documentclass[11pt,oneside,openany,report]{jsbook}

\usepackage[a4paper,truedimen,margin=25truemm]{geometry}
\usepackage{cscover}
\usepackage[dvipdfmx]{graphicx}
\usepackage[nobreak]{cite}
\usepackage[a4paper,dvipdfmx,pdfdisplaydoctitle=true,%
    bookmarks=true,bookmarksnumbered=true,bookmarkstype=toc,bookmarksopen=true,%
    pdftitle={周辺視野を含むコントラスト感度関数の測定とモデル化に関する研究},%
    pdfauthor={桑木啓伍}%
    ]{hyperref}
\usepackage{pxjahyper}

\renewcommand{\bibname}{参考文献}
\setcounter{tocdepth}{2}
\pagestyle{plain}

\newcommand{\TODO}[1]{\textbf{[TODO: #1]}}
%\renewcommand{\TODO}[1]{}

\thesistype{学士特定課題研究論文}
\title{周辺視野を含むコントラスト感度関数の測定とモデル化に関する研究}
\author{桑木 啓伍}
\studentid{23B30351}
\affiliation{東京科学大学\\情報理工学院\\情報工学系} 
\date{2026年1月}

\supervisorname{指導教員}
\supervisor{齋藤 豪}
%\dsupervisorname{副指導教員}
%\dsupervisor{工学 次郎}

\begin{document}

\frontmatter
\maketitle

\chapter{概要}

\tableofcontents
\listoffigures
\listoftables

%%%%%%%%%%%%%%%%%%%%%%%%%%%%%%%%%%%%%%%%%%%%%%%%%%

\mainmatter
\chapter{序論}
本論文は、学士特定課題研究論文の書き方\cite{tokodai-xyz2015}の一例を示す。

\section{背景と目的}

%%%%%%%%%%%%%%%%%%%%%%%%%%%%%%%%%%%%%%%%%%%%%%%%%%

\chapter{関連研究}

  \section{はじめに}

  \section{視覚系の基礎知識}

    \subsection{視覚情報処理}

    \subsection{輝度と明るさ}

    \subsection{色空間}

    \subsection{色収差}

  \section{CSFに関する基礎知識}

    \subsection{コントラスト}

    \subsection{CSF}

  \section{視覚実験手法}

    \subsection{調整法}

    \subsection{極限法}

    \subsection{恒常法}

    \subsection{Ψ 法}

  \section{おわりに}

%%%%%%%%%%%%%%%%%%%%%%%%%%%%%%%%%%%%%%%%%%%%%%%%%%

\chapter{輝度および赤緑コントラスト感度関数の測定実験}

  \section{はじめに}

  \section{輝度コントラスト感度関数測定実験}

    \subsection{実験手法}

    \subsection{測定結果}

    \subsection{考察}

  \section{赤緑コントラスト感度関数測定実験}

    \subsection{実験手法}

    \subsection{測定結果}

    \subsection{考察}

  \section{おわりに}

%%%%%%%%%%%%%%%%%%%%%%%%%%%%%%%%%%%%%%%%%%%%%%%%%%

\chapter{コントラスト感度関数モデルへの適用}

  \section{はじめに}

  \section{おわりに}

%%%%%%%%%%%%%%%%%%%%%%%%%%%%%%%%%%%%%%%%%%%%%%%%%%

\chapter{結論}

  \section{本研究のまとめ}

  \section{今後の課題}

    \subsection{コントラスト感度関数測定実験}

    \subsection{コントラスト感度関数モデル}

%%%%%%%%%%%%%%%%%%%%%%%%%%%%%%%%%%%%%%%%%%%%%%%%%%

\appendix
\chapter{CSF測定実験結果の詳細データ}

%%%%%%%%%%%%%%%%%%%%%%%%%%%%%%%%%%%%%%%%%%%%%%%%%%

\backmatter
\chapter{謝辞}

%%%%%%%%%%%%%%%%%%%%%%%%%%%%%%%%%%%%%%%%%%%%%%%%%%

\bibliographystyle{jplain}
\bibliography{references}
%\begin{thebibliography}{99}
%  \bibitem{tokodai-xyz2015} 科学大太郎. 良い論文の書き方. \textit{Journal of XYZ}, Vol.~3, No.~4, pp. 15--34, 2015.
%\end{thebibliography}

\end{document}
